\chapter{先辨认问题：通信模型的优化地图}\label{ch:models}
\chaptertakeaway{把系统需求翻译成变量、目标和约束，并从数学结构判断求解难度。后续各章只改变模型中的少数部件，使同一条思路能迁移到新系统。}

\section{优化的对象是什么}
通信系统的优化通常是在有限功率、带宽、天线和时隙下，选择可实施的收发策略。优化并不自动决定什么是好的系统：总吞吐量、最差用户速率、能效和时延对应不同设计。必须先说明服务需求，再决定目标。

设 $x$ 汇集设计变量，$\xi$ 表示信道、到达流量等环境量。一次确定性设计可写为
\begin{equation}\label{eq:general}
 \min_{x\in\mathcal X} f_0(x;\xi),\qquad
 f_i(x;\xi)\le0\ (i=1,\ldots,m),\quad
 h_j(x;\xi)=0\ (j=1,\ldots,r).
\end{equation}
集合 $\mathcal X$ 规定变量本身的限制，例如非负功率、复数波束或二元指派。函数 $f_i$ 表示功率预算、最低服务质量等约束。最大化效用 $U$ 可以写成最小化 $-U$，但全文推导时会保留最自然的方向。

建模先完成四个具体选择。第一，决定哪些量由设计者控制：发射功率是变量，给定时刻测得的信道一般是参数。第二，确定决策发生在信道观测之前还是之后；这决定它是固定策略还是随信道变化的策略。第三，说明约束是每时隙成立还是仅在长期平均下成立。第四，确定物理层速率模型，例如是否允许高斯码本、将干扰作为噪声、是否忽略有限码长损失。换一个选择，通常就换了一个数学问题。

\section{用一个速率公式连接主要变量}
先考虑 $K$ 条单天线链路，每条链路每次发送一个独立数据流。令 $p_k\ge0$ 为第 $k$ 个发射机的功率，$g_{kj}\ge0$ 为发射机 $j$ 到接收机 $k$ 的功率增益，$\sigma_k^2>0$ 为接收噪声功率。将其他数据流作为噪声时，
\begin{equation}\label{eq:sinr-basic}
 \gamma_k(p)=\frac{g_{kk}p_k}{\sigma_k^2+\sum_{j\ne k}g_{kj}p_j},\qquad
 R_k(p)=B_k\log_2(1+\gamma_k(p)).
\end{equation}
$B_k$ 是带宽，单位 Hz，$R_k$ 的单位是 bit/s；$\log_2(1+\gamma_k)$ 是频谱效率，单位 bit/s/Hz。功率与噪声必须采用同一线性单位。以 dBm 给出的量要先转成瓦特，不能直接代入加法。

当用户正交传输时，交叉增益项消失，$R_k$ 对 $p_k$ 是凹函数；当各用户共享频谱时，$p_j$ 同时增加用户 $j$ 的信号和用户 $k$ 的干扰，联合曲率发生变化。由此出现本书最重要的分界：\emph{通信公式中的对数并不会自动使整个问题凸化。}

若设计带宽 $b_k$，噪声功率也随带宽改变。对平坦信道，令 $a_k=|h_k|^2/N_0$，其中 $N_0$ 是采用本模型约定的噪声功率谱密度，单位 W/Hz，则
\begin{equation}\label{eq:band-rate}
 R_k(b_k,p_k)=b_k\log_2\left(1+\frac{a_kp_k}{b_k}\right),\quad b_k>0.
\end{equation}
此时不能把噪声总功率固定后只在速率前乘带宽。式\eqref{eq:band-rate} 是凹函数的透视形式，将导向凸的联合资源分配。

\section{通信问题与数学形式的对应}
表\ref{tab:map} 先按变量与耦合方式分类。同一个系统可能同时落入数行；例如带用户选择的稳健波束设计同时包含离散变量、不确定性和二次约束。

\begin{longtable}{p{.19\textwidth}p{.27\textwidth}p{.28\textwidth}p{.15\textwidth}}
\caption{常见通信模型的结构地图。表中的凸性依赖所列条件。}\label{tab:map}\\
\toprule 系统任务 & 变量和典型目标 & 抽象结构与关键条件 & 主要章节\\\midrule\endfirsthead
\toprule 系统任务 & 变量和典型目标 & 抽象结构与关键条件 & 主要章节\\\midrule\endhead
正交资源分配 & 功率/带宽；和速率最大 & 凹效用最大化；透视函数 & 第3章\\
固定服务质量功控 & 功率；总功率最小 & 固定 SINR 门限可化 LP，也可写 GP & 第3章\\
最差用户公平 & 功率；最小 SINR 最大 & 超水平集可行性为 LP；二分 & 第3章\\
多天线单流下行 & 复波束；总功率最小 & 单播 SINR 约束可化 SOCP & 第4章\\
多播或单位模设计 & 波束/相位；功率或速率 & 非凸 QCQP；SDR、SCA、交替法 & 第4、6章\\
信道估计/检测 & 信道向量或符号 & LS、正则化；离散检测则为整数问题 & 第4、5、7章\\
干扰网络吞吐量 & 多用户功率/波束 & 非凸和速率；WMMSE、SCA & 第6章\\
能效优化 & 功率；速率/总耗电 & 分式规划；凹分子时有全局方法 & 第6章\\
用户调度与指派 & 二元变量；固定收益最大 & 指派 LP 可整数；一般为混合整数 & 第7章\\
网络路由与拥塞 & 链路流量；效用最大 & 流守恒 LP、凸网络效用、对偶分解 & 本章、第5章\\
不确定信道设计 & 策略；最坏或平均性能 & 稳健、机会约束、随机优化 & 第4、8章\\
时延与在线控制 & 每时隙动作；长期效用 & 随机过程、队列约束 & 第8章\\\bottomrule
\end{longtable}

\subsection{公平性改变目标，也可能改变求解方法}
设 $R_k>0$。加权和速率 $\sum_k w_kR_k$ 将资源倾向于边际收益大的用户，$w_k>0$ 是设计权重。比例公平常用 $\sum_k\log R_k$，它对低速率用户更敏感。最差用户速率 $\min_k R_k$ 则直接保护瓶颈。通过新增变量 $t$，后者写为
\begin{equation}
 \max_{x,t}t,\qquad R_k(x)\ge t\quad(\forall k).
\end{equation}
若 $R_k$ 凹，这些约束是凸约束；若是干扰网络中的式\eqref{eq:sinr-basic}，仍需进一步利用固定门限的结构。公平性是服务政策，不能只凭算法容易程度来选择。

\subsection{网络层的耦合通常来自容量与流守恒}
考虑有向网络，$x_e^{(s)}\ge0$ 为业务 $s$ 在链路 $e$ 的流量，$r_s\ge0$ 为其端到端速率。对每个中间节点，流入总量等于流出总量；在源、目的节点分别相差 $r_s$。若链路容量 $c_e$ 固定，则
\begin{equation}
 \max \sum_s U_s(r_s),\qquad
 \sum_s x_e^{(s)}\le c_e,\quad\text{各业务满足流守恒}.
\end{equation}
当 $U_s$ 凹时，这是凸优化。若容量由共享频谱的干扰功率控制决定，约束右侧变成非凹的 $c_e(p)$，原有凸性通常丢失。对偶分解把容量约束的乘子解释成链路拥塞价格，见第5章。

\subsection{估计和检测共用观测模型，却有不同难度}
线性观测写为 $y=Ax+n$：$y$ 是已知接收样本，$A$ 是导频或信道矩阵，$x$ 是待估计量。连续 $x$ 的最小二乘 $\min_x\norm{Ax-y}_2^2$ 是凸二次问题；加入 $\ell_1$ 正则化仍凸。若 $x$ 必须属于 QAM 星座的笛卡尔积，则是离散最小二乘，不能沿用连续问题的全局保证。连续松弛仍有价值，因为它提供候选解与可认证的界。

\section{从一个模型到一份可验证的答案}
图\ref{fig:workflow} 给出解题顺序。先检查空可行集、零预算、信道增益为零等边界情况，再识别结构，能避免在不可行模型上反复调参。
\begin{figure}[htbp]\centering
\begin{tikzpicture}[node distance=7mm,box/.style={draw,rounded corners,align=center,text width=10.5cm,inner sep=7pt},>=Stealth]
\node[box] (a) {系统假设、CSI 可得时间、变量、单位与服务要求};
\node[box,below=of a] (b) {写出目标与约束；检查可行性和尺度};
\node[box,below=of b] (c) {识别凸性、单调性、可分性、整数性与不确定性};
\node[box,below=of c] (d) {优先精确变换；其次用有界松弛或条件明确的局部算法};
\node[box,below=of d] (e) {回到原模型核验残差、目标值与最优性证书};
\draw[->] (a)--(b);\draw[->] (b)--(c);\draw[->] (c)--(d);\draw[->] (d)--(e);
\end{tikzpicture}
\caption{模型到验证的流程。每个箭头代表需完成的判断；求解器状态只覆盖最后一步的一部分。}\label{fig:workflow}
\end{figure}

\emph{精确变换}保留可行解与目标关系，例如相位固定后的单播 SOCP；\emph{松弛}扩大可行集，例如去掉秩一约束；\emph{近似}改变函数或可行域，例如在当前点构造凹下界。三者必须分清：它们分别支持等价最优、目标界或局部迭代等不同结论。

对最小化问题，报告可行解给出上界，松弛或对偶解给出下界；对最大化方向正好相反。若上下界之差小，可以量化接近最优的程度。只报告“算法收敛”通常指迭代变化变小，还需要说明收敛到的是什么。
